\makeatletter
\def\input@path{{./}{../}{../../}{../../../}{../../../../}{preamble/}{../preamble/}{../../preamble/}{../../../preamble/}}
\makeatother
\input{preamble_exams.tex}
\input{graphs/highlights.tex}
% \GTenHighlightsfalse % Uncomment to disable instructional highlights.
\input{graphs/specific_highlights.tex}
\CoordinatedTermTwoPalette
\usepackage{needspace}
\usepackage{titling}
\setlength{\droptitle}{-2em}
\pretitle{\par\begingroup\centering\LARGE}
\posttitle{\par\endgroup}
\preauthor{}\postauthor{}
\predate{}\postdate{}
\maketitlehookd{\vspace{-1em}}
\linespread{1.12}
\setlength{\parskip}{5pt}
\newgeometry{margin=0.65in,top=0.55in,bottom=0.65in}
% Local table presentation reuses the coordinated registry's semantic roles.
% No shared palette or registry entry is changed by this activity.
\newcommand{\ExampleOne}{\textbullet\ \hyperref[ex:intro-one]{Example 1}: }
\newcommand{\ExampleTwo}{\textbullet\ \hyperref[ex:intro-two]{Example 2}: }
\title{Grade 11 Term 2 Home Practice -- C4\\\large Confidence intervals: construct and recover\\[0.25em]
\normalsize Global Economics -- 2026--2027}
\author{}
\date{}
\begin{document}
\maketitle
\setcounter{tocdepth}{1}
\setcounter{secnumdepth}{2}
\phantomsection\label{toc}
\tableofcontents
\section*{Evaluated criteria}

\begin{description}
\item[\textbf{C4}] Constructs confidence intervals for a population mean with known variance using the normal distribution, applying them to contexts such as average income, average investment return, or average price level.
\end{description}
\textbf{How to practise.} Try each task before opening its blue solution link. Show the model, algebra, units and interpretation. The arrows return to the contents, current problem or problem list. The inverse tasks help you check how the C4 interval is constructed.

\clearpage
\section{Introduction}\label{sec:introduction}
\subsection*{Highlight legend}
\phantomsection\label{tab:highlight-legend}
\begin{tabularx}{\linewidth}{@{}>{\raggedright\arraybackslash}p{.23\linewidth}>{\raggedright\arraybackslash}p{.25\linewidth}>{\raggedright\arraybackslash}X@{}}
\toprule
\textbf{Symbol or notation}&\textbf{Concept}&\textbf{Brief meaning}\\\midrule
$\HLPopulationMean{\mu}$&\HLPopulationMean{Population mean}&Fixed mean to be estimated.\\
$\HLSampleMean{\bar x}$&\HLSampleMean{Sample mean}&Mean of the sampled values.\\
$\HLPopulationSD{\sigma}$&\HLPopulationSD{Population SD}&Known spread of individual values.\\
$\HLObservationCount{n}$&\HLObservationCount{Sample size}&Number of sampled units.\\
$\HLConfidence{C}$&\HLConfidence{Confidence level}&Long-run coverage of the procedure.\\
$\HLCriticalValue{z^*}$&\HLCriticalValue{Critical value}&Positive standard normal cutoff.\\
$\HLStandardError{\mathrm{SE}}$&\HLStandardError{Standard error}&Spread of the sample mean.\\
$\HLMarginOfError{E}$&\HLMarginOfError{Margin of error}&Distance from center to endpoint.\\
$[\HLIntervalBounds{L},\HLIntervalBounds{U}]$&\HLIntervalBounds{Interval bounds}&Lower and upper endpoints.\\
\bottomrule
\end{tabularx}

Names and units identify each role with or without color. Related shades connect related ideas; definitions distinguish them. See the \hyperref[tab:notation-reference]{detailed notation reference} for contextual meanings.

\subsection{Introductory example 1: a courier earnings estimate}\label{ex:intro-one}
A cooperative estimates the \HLPopulationMean{population mean} $\HLPopulationMean{\mu}$ of reference-week earnings for its 5,000 couriers. A simple random sample of $\HLObservationCount{100}$ distinct couriers has \HLSampleMean{sample mean} $\HLSampleMean{\bar x=420}$ thousand COP. Prior population records establish $\HLPopulationSD{\sigma=50}$ thousand COP; assume no extreme skew or outliers. The sampling fraction is 2\%, so approximate independence and an approximately normal sample mean are reasonable.

The \HLStandardError{standard error} is $\HLStandardError{\mathrm{SE}}=\HLPopulationSD{50}/\sqrt{\HLObservationCount{100}}=\HLStandardError{5}$ thousand COP. At $\HLConfidence{C=0.95}$, using $\HLCriticalValue{z^*=1.96}$ gives $\HLMarginOfError{E=9.8}$ thousand COP. Thus the estimate is centered at 420 with endpoints 410.2 and 429.8 thousand COP. It estimates a mean, not where 95\% of individual couriers' earnings lie.

\subsection{Introductory example 2: reading a warehouse report}\label{ex:intro-two}
A warehouse report gives a $\HLConfidence{90\%}$ interval $[\HLIntervalBounds{48.355},\HLIntervalBounds{51.645}]$ USD per pallet for the mean charge across 4,000 pallets. It used the known-population-$\sigma$ $z$ procedure with construction cutoff $\HLCriticalValue{1.645}$ and $\HLObservationCount{n=100}$. Receipts were sampled randomly, one per distinct pallet; the 2.5\% sampling fraction and approximately bell-shaped charges support the model.

The midpoint is $\HLSampleMean{50}$ USD and the half-width is $\HLMarginOfError{1.645}$ USD. Reading an interval backward can reveal its center and the spread used in its construction. \textbf{Quick check:} does a USD 52 charge on one pallet contradict this report? No: an individual charge is not the population mean.

\clearpage
\section{Notation Reference and Practice Procedure}\label{sec:notation-and-procedure}
\subsection{Notation reference}\label{tab:notation-reference}
\begingroup\small
\setlength{\tabcolsep}{4pt}
\renewcommand{\arraystretch}{1.15}
\begin{longtable}{@{}>{\raggedright\arraybackslash}p{.15\linewidth}>{\raggedright\arraybackslash}p{.18\linewidth}>{\raggedright\arraybackslash}p{.27\linewidth}>{\raggedright\arraybackslash}p{\dimexpr.40\linewidth-6\tabcolsep\relax}@{}}
\toprule
\textbf{Symbol or notation}&\textbf{Concept}&\textbf{Definition}&\textbf{Meaning in the introductory examples}\\\midrule\endfirsthead
\toprule
\textbf{Symbol or notation}&\textbf{Concept}&\textbf{Definition}&\textbf{Meaning in the introductory examples}\\\midrule\endhead
$\HLPopulationMean{\mu}$&\HLPopulationMean{Population mean}&Fixed arithmetic mean of the defined population.&\ExampleOne Mean reference-week earnings of all 5,000 couriers.\par\ExampleTwo Mean charge across all 4,000 pallets.\\
$\HLSampleMean{\bar x}$&\HLSampleMean{Sample mean}&Arithmetic mean of the sample; center of this symmetric interval.&\ExampleOne $\HLSampleMean{420}$ thousand COP.\par\ExampleTwo Report's midpoint: $\HLSampleMean{50}$ USD.\\
$\HLPopulationSD{\sigma}$; $s$&\HLPopulationSD{Population SD}; sample SD&$\sigma$ describes population spread; $s$ estimates spread from sample data. This procedure requires known $\sigma$, not just $s$.&\ExampleOne Known $\HLPopulationSD{50}$ thousand COP.\par\ExampleTwo Population spread used by the report, recoverable from its construction.\\
$\sigma^2$&Population variance&Square of population SD; units are squared.&\ExampleOne $50^2=2{,}500$ in squared thousand COP.\\
$\HLObservationCount{n}$; $N$&\HLObservationCount{Sample size}; population size&$n$ counts sampled units; $N$ counts population units. Without replacement, $n/N\leq 0.10$ supports approximate independence.&\ExampleOne $\HLObservationCount{100}$ of 5,000 couriers.\par\ExampleTwo $\HLObservationCount{100}$ of 4,000 pallets.\\
$\HLConfidence{C}$&\HLConfidence{Confidence level}&Decimal long-run proportion of intervals covering $\mu$ under repeated sampling.&\ExampleOne $\HLConfidence{0.95}$ (95\%).\par\ExampleTwo $\HLConfidence{0.90}$ (90\%).\\
$\Phi$; $\Phi^{-1}$&Normal CDF; inverse CDF&$\Phi(z)=P(Z\leq z)$ for $Z\sim N(0,1)$; the inverse converts cumulative area to a cutoff.&\ExampleOne Left area $(1+0.95)/2=0.975$.\par\ExampleTwo Left area $0.95$.\\
$\HLCriticalValue{z^*}$&\HLCriticalValue{Critical value}&Positive cutoff leaving central area $C$ and $(1-C)/2$ in each tail.&\ExampleOne Construction cutoff $\HLCriticalValue{1.96}$.\par\ExampleTwo Construction cutoff $\HLCriticalValue{1.645}$.\\
$\HLStandardError{\mathrm{SE}}$&\HLStandardError{Standard error}&$\mathrm{SE}=\sigma/\sqrt n$, the SD of the sample mean under the model.&\ExampleOne $\HLStandardError{5}$ thousand COP.\\
$\HLMarginOfError{E}$&\HLMarginOfError{Margin of error}&$E=z^*\sigma/\sqrt n$; half the interval width, in the response's units.&\ExampleOne $\HLMarginOfError{9.8}$ thousand COP.\par\ExampleTwo $\HLMarginOfError{1.645}$ USD.\\
$\HLIntervalBounds{L}$, $\HLIntervalBounds{U}$&\HLIntervalBounds{Interval bounds}&$L=\bar x-E$ and $U=\bar x+E$; total width is $U-L=2E$.&\ExampleOne 410.2 and 429.8 thousand COP.\par\ExampleTwo $\HLIntervalBounds{48.355}$ and $\HLIntervalBounds{51.645}$ USD.\\
\bottomrule
\end{longtable}
\endgroup

\clearpage
\subsection{Relationships and practice procedure}
For a two-sided interval, write the confidence level as a decimal, for example $\HLConfidence{95\%=0.95}$. The theoretical cutoff is
\[
\HLCriticalValue{z^*}=\Phi^{-1}\!\left(\frac{1+\HLConfidence{C}}{2}\right).
\]
Common numerical approximations are 1.645 (90\%), 1.960 (95\%) and 2.576 (99\%). In inverse problems, use the \emph{stated construction cutoff}, even if a calculator gives more digits. The supplied endpoints are exact relative to that cutoff unless explicitly described as rounded.
\[
\HLMarginOfError{E}=\frac{\HLCriticalValue{z^*}\HLPopulationSD{\sigma}}{\sqrt{\HLObservationCount{n}}},
\qquad [\HLIntervalBounds{L},\HLIntervalBounds{U}]
=[\HLSampleMean{\bar x}-\HLMarginOfError{E},\HLSampleMean{\bar x}+\HLMarginOfError{E}].
\]
\[
\HLSampleMean{\bar x}=\frac{\HLIntervalBounds{L}+\HLIntervalBounds{U}}{2},
\qquad \HLMarginOfError{E}=\frac{\HLIntervalBounds{U}-\HLIntervalBounds{L}}{2}.
\]
Starting from $E\sqrt n=z^*\sigma$ gives the inverse rules
\[
\HLObservationCount{n}=\left(\frac{\HLCriticalValue{z^*}\HLPopulationSD{\sigma}}{\HLMarginOfError{E}}\right)^2,
\qquad \HLPopulationSD{\sigma}=\frac{\HLMarginOfError{E}\sqrt{\HLObservationCount{n}}}{\HLCriticalValue{z^*}}.
\]
\begin{enumerate}
\item \textbf{Define the model.} Identify the population, its mean $\mu$ and the response units. Check random sampling, approximate independence and a normal or sufficiently well-behaved large-sample mean. Confirm that $\sigma$ is known; $s$ alone does not justify this known-$\sigma$ procedure.
\item \textbf{Organize the three stages.} For each general problem, construct the full $2\times2$ grid of sample sizes and confidence levels while keeping $\bar x$ and $\sigma$ fixed. Then solve each of the four inverse reports independently; never combine reports as simultaneous equations for shared unknowns.
\item \textbf{List knowns and unknowns.} Given $\bar x,\sigma,n,C$, construct $[L,U]$. Given $[L,U],C,\sigma$, recover $\bar x,n$. Given $[L,U],C,n$, recover $\bar x,\sigma$ only when the report explicitly used the known-$\sigma$ $z$ procedure.
\item \textbf{Choose the cutoff and solve.} Use the normal inverse for a forward problem or the supplied construction cutoff for an inverse problem. Find the midpoint and half-width first when endpoints are given. Rearrange symbolically before substituting; retain precision until the final report.
\item \textbf{Check.} Reconstruct the original endpoints. $n$ is a positive integer and $\sigma,E$ are positive. Rounded endpoints may not uniquely identify the original $n$: do not automatically round upward. Upward rounding belongs to a different question, planning a \emph{minimum} sample size to meet a margin target.
\item \textbf{Interpret and compare.} The interval estimates $\mu$, not individual observations. At fixed other inputs, larger $C$ increases $z^*$ and width; larger $n$ reduces width by $1/\sqrt n$; larger $\sigma$ increases width proportionally.
\end{enumerate}
\textbf{Confidence meaning.} We use ``95\% confident'' to describe a procedure that captures the fixed population mean in about 95\% of repeated random samples under its assumptions. It is not a 95\% probability assigned to the fixed $\mu$ after observing this particular interval.
\secInicio[][sec:notation-and-procedure]

\clearpage
\section{General Worked Problems}\label{sec:general-worked}
\phantomsection\label{sec:worked-solution-examples}
\SectionProblemsTableOfContents{
\SectionProblemLink{prob:c4-1}{Courier earnings: construct and recover (C4)}\\
\SectionProblemLink{prob:c4-2}{Warehouse charges: construct and recover (C4)}\\
\SectionProblemLink{prob:c4-3}{Freight charges: construct and recover (C4)}
}
\par\vspace{5mm}
Each problem has three stages: four forward constructions, two separate reports to recover
$n$ and $\bar x$, and two further separate reports to recover $\bar x$ and $\sigma$.
Treat every inverse report independently: its unknowns are not shared with another report
or with Stage 1. All inverse endpoints are exact relative to the supplied construction cutoff.

\subsection{Courier earnings: construct and recover}\label{prob:c4-1}
A courier cooperative studies reference-week earnings across its 4,000 members.
In every case, use a simple random sample of distinct couriers without replacement.
Assume no extreme skew or outliers and an approximately normal sample mean for all the stated
or recovered sizes. The sample fractions are below 10\%, supporting approximate independence;
use the stated known-$\sigma$ $z$ model without a finite-population correction.
The population size is unchanged across the separate scenarios; report-specific means and SDs may differ.
All response values, including SDs and interval endpoints, are in thousand COP.
\QuestionsTasks[sec:general-worked]{
{4}/{sol:c4-1-a}/{\textbf{Stage 1 -- Construct four intervals.}
Use $\bar x=410$ and known $\sigma=80$ thousand COP throughout.
For $n_1=64$ and $n_2=256$, construct all four two-sided intervals at
$C_1=95\%$ ($z^*=1.96$) and $C_2=99\%$ ($z^*=2.576$).
Show the four combinations in a comparison table, including margins and widths.
Interpret the estimates and explain both width effects.},
{4}/{sol:c4-1-b}/{\textbf{Stage 2 -- Recover $n$ and $\bar x$.} For EACH independent report below, recover both the sample size and the sample mean. The known-population-$\sigma$ $z$ procedure was used; $n$ and $\bar x$ are omitted.
\newline Report 1: interval $[390.4,429.6]$ thousand COP; $C=95\%$, construction cutoff $z^*=1.96$, known $\sigma=80$ thousand COP.
\newline Report 2: interval $[404.392,445.608]$ thousand COP; $C=99\%$, construction cutoff $z^*=2.576$, known $\sigma=80$ thousand COP.
\newline Find each midpoint and half-width first; show inverse algebra, reconstruct the endpoints and interpret each interval.},
{4}/{sol:c4-1-c}/{\textbf{Stage 3 -- Recover $\bar x$ and $\sigma$.} For EACH further independent report below, recover both the sample mean and the population SD. Each report was constructed using the \textbf{known-population-$\sigma$ $z$ procedure}; the population SD used by the analyst was omitted from the report, as was $\bar x$.
\newline Report 3: interval $[416.84,443.16]$ thousand COP; $C=90\%$, construction cutoff $z^*=1.645$, sample size $n=100$.
\newline Report 4: interval $[405.2,424.8]$ thousand COP; $C=95\%$, construction cutoff $z^*=1.96$, sample size $n=144$.
\newline Find each midpoint and half-width first; show inverse algebra, reconstruct the endpoints and interpret the recovered quantities. Explain why the recovered SD is $\sigma$, not $s$.}
}
\clearpage
\subsubsection*{C4 -- Stage 1: four interval constructions}\label{sol:c4-1-a}
The parameter $\HLPopulationMean{\mu}$ is the cooperative's mean reference-week earnings, in thousand COP. The sampling and shape assumptions justify the approximate normal model. The known \HLPopulationSD{population SD} is $\HLPopulationSD{\sigma=80}$ thousand COP; the fixed \HLSampleMean{sample mean} is $\HLSampleMean{\bar x=410}$ thousand COP.
For each supplied confidence level, $z^*\approx\Phi^{-1}((1+C)/2)$.
Use the supplied cutoffs in the numerical calculations. The formulas are
\[
\HLStandardError{\mathrm{SE}}=\frac{\HLPopulationSD{\sigma}}{\sqrt{\HLObservationCount{n}}},\quad
\HLMarginOfError{E}=\HLCriticalValue{z^*}\HLStandardError{\mathrm{SE}},\quad
[\HLIntervalBounds{L},\HLIntervalBounds{U}]=[\HLSampleMean{\bar x}-\HLMarginOfError{E},\HLSampleMean{\bar x}+\HLMarginOfError{E}],\quad W=2E.
\]
\textbf{$n=64$, $C=95\%$.}
\[\begin{aligned}
\HLStandardError{\mathrm{SE}}&=\frac{\HLPopulationSD{80}}{\sqrt{\HLObservationCount{64}}}=\frac{80}{8}=\HLStandardError{10},\qquad \HLMarginOfError{E}=\HLCriticalValue{1.96}(\HLStandardError{10})=\HLMarginOfError{19.6},\\
[L,U]&=[\HLSampleMean{410}-\HLMarginOfError{19.6},\HLSampleMean{410}+\HLMarginOfError{19.6}]=[\HLIntervalBounds{390.4},\HLIntervalBounds{429.6}],\qquad W=2(19.6)=39.2.
\end{aligned}\]
\textbf{$n=64$, $C=99\%$.}
\[\begin{aligned}
\HLStandardError{\mathrm{SE}}&=\frac{\HLPopulationSD{80}}{\sqrt{\HLObservationCount{64}}}=\frac{80}{8}=\HLStandardError{10},\qquad \HLMarginOfError{E}=\HLCriticalValue{2.576}(\HLStandardError{10})=\HLMarginOfError{25.76},\\
[L,U]&=[\HLSampleMean{410}-\HLMarginOfError{25.76},\HLSampleMean{410}+\HLMarginOfError{25.76}]=[\HLIntervalBounds{384.24},\HLIntervalBounds{435.76}],\qquad W=2(25.76)=51.52.
\end{aligned}\]
\textbf{$n=256$, $C=95\%$.}
\[\begin{aligned}
\HLStandardError{\mathrm{SE}}&=\frac{\HLPopulationSD{80}}{\sqrt{\HLObservationCount{256}}}=\frac{80}{16}=\HLStandardError{5},\qquad \HLMarginOfError{E}=\HLCriticalValue{1.96}(\HLStandardError{5})=\HLMarginOfError{9.8},\\
[L,U]&=[\HLSampleMean{410}-\HLMarginOfError{9.8},\HLSampleMean{410}+\HLMarginOfError{9.8}]=[\HLIntervalBounds{400.2},\HLIntervalBounds{419.8}],\qquad W=2(9.8)=19.6.
\end{aligned}\]
\textbf{$n=256$, $C=99\%$.}
\[\begin{aligned}
\HLStandardError{\mathrm{SE}}&=\frac{\HLPopulationSD{80}}{\sqrt{\HLObservationCount{256}}}=\frac{80}{16}=\HLStandardError{5},\qquad \HLMarginOfError{E}=\HLCriticalValue{2.576}(\HLStandardError{5})=\HLMarginOfError{12.88},\\
[L,U]&=[\HLSampleMean{410}-\HLMarginOfError{12.88},\HLSampleMean{410}+\HLMarginOfError{12.88}]=[\HLIntervalBounds{397.12},\HLIntervalBounds{422.88}],\qquad W=2(12.88)=25.76.
\end{aligned}\]
All SEs, margins, endpoints and widths above and below are in thousand COP.
\begin{center}
\renewcommand{\arraystretch}{1.25}
\begin{tabular}{@{}rrrrll@{}}\toprule
$\HLObservationCount{n}$&$\HLConfidence{C}$&$\HLCriticalValue{z^*}$&$\HLMarginOfError{E}$&$[\HLIntervalBounds{L},\HLIntervalBounds{U}]$&Width $W$\\\midrule
$\HLObservationCount{64}$&$\HLConfidence{95\%}$&$\HLCriticalValue{1.96}$&$\HLMarginOfError{19.6}$&$[\HLIntervalBounds{390.4},\HLIntervalBounds{429.6}]$&$39.2$\\
$\HLObservationCount{64}$&$\HLConfidence{99\%}$&$\HLCriticalValue{2.576}$&$\HLMarginOfError{25.76}$&$[\HLIntervalBounds{384.24},\HLIntervalBounds{435.76}]$&$51.52$\\
$\HLObservationCount{256}$&$\HLConfidence{95\%}$&$\HLCriticalValue{1.96}$&$\HLMarginOfError{9.8}$&$[\HLIntervalBounds{400.2},\HLIntervalBounds{419.8}]$&$19.6$\\
$\HLObservationCount{256}$&$\HLConfidence{99\%}$&$\HLCriticalValue{2.576}$&$\HLMarginOfError{12.88}$&$[\HLIntervalBounds{397.12},\HLIntervalBounds{422.88}]$&$25.76$\\
\bottomrule\end{tabular}\end{center}
At each row's confidence level, we are 95\% or 99\% confident, respectively, that the cooperative's mean reference-week earnings lies within that row's bounds.
These intervals estimate a population mean, not a range of individual values.
Holding $\bar x$ and $\sigma$ fixed, increasing $n$ from 64 to 256 multiplies width by
$\sqrt{64/256}=1/2$: quadrupling the sample size halves width at either confidence level.
Increasing confidence from 95\% to 99\% increases width by the factor
$2.576/1.96\approx 1.3143$ at either sample size, since the cutoff increases.
Every interval remains centered at 410 thousand COP. These are controlled comparisons;
an actual new random sample could have a different sample mean.
\secInicio[prob:c4-1][sec:general-worked]
\clearpage
\subsubsection*{C4 -- Stage 2: recover $n$ and $\bar x$}\label{sol:c4-1-b}
Each report estimates the cooperative's mean reference-week earnings. Work independently, using the cutoff stated in that report. All response calculations on this page are in thousand COP; $n$ is a count of couriers.
\par\noindent\textbf{Report 1: midpoint and half-width first.}
\[\begin{aligned}
\HLSampleMean{\bar x}&=\frac{L+U}{2}=\frac{390.4+429.6}{2}=\frac{820}{2}=\HLSampleMean{410},\\
\HLMarginOfError{E}&=\frac{U-L}{2}=\frac{429.6-390.4}{2}=\frac{39.2}{2}=\HLMarginOfError{19.6}.
\end{aligned}\]
\[E=\frac{z^*\sigma}{\sqrt n}\ \Longrightarrow\ E\sqrt n=z^*\sigma\ \Longrightarrow\ \sqrt n=\frac{z^*\sigma}{E}\ \Longrightarrow\ n=\left(\frac{z^*\sigma}{E}\right)^2.\]
\[\HLObservationCount{n}=\left(\frac{\HLCriticalValue{1.96}\cdot\HLPopulationSD{80}}{\HLMarginOfError{19.6}}\right)^2=\left(\frac{156.8}{19.6}\right)^2=8^2=\HLObservationCount{64}\text{ couriers}.\]
The response units cancel in the ratio. Thus $\bar x=410$ thousand COP and $n=64$ exactly; no upward rounding is used.
\textbf{Endpoint verification.} The recovered sampling fraction is $64/4000\approx 1.60\%<10\%$. Recompute
\[\begin{aligned}
\HLStandardError{\mathrm{SE}}&=\frac{80}{\sqrt{64}}=\frac{80}{8}=\HLStandardError{10},\qquad \HLMarginOfError{E}=1.96(10)=\HLMarginOfError{19.6},\\
[\bar x-E,\bar x+E]&=[410-19.6,410+19.6]=[\HLIntervalBounds{390.4},\HLIntervalBounds{429.6}].
\end{aligned}\]
This reproduces the supplied endpoints exactly relative to $z^*=1.96$. We are 95\% confident that the cooperative's mean reference-week earnings is between 390.4 and 429.6 thousand COP.
\par\noindent\textbf{Report 2: midpoint and half-width first.}
\[\begin{aligned}
\HLSampleMean{\bar x}&=\frac{L+U}{2}=\frac{404.392+445.608}{2}=\frac{850}{2}=\HLSampleMean{425},\\
\HLMarginOfError{E}&=\frac{U-L}{2}=\frac{445.608-404.392}{2}=\frac{41.216}{2}=\HLMarginOfError{20.608}.
\end{aligned}\]
\[E=\frac{z^*\sigma}{\sqrt n}\ \Longrightarrow\ E\sqrt n=z^*\sigma\ \Longrightarrow\ \sqrt n=\frac{z^*\sigma}{E}\ \Longrightarrow\ n=\left(\frac{z^*\sigma}{E}\right)^2.\]
\[\HLObservationCount{n}=\left(\frac{\HLCriticalValue{2.576}\cdot\HLPopulationSD{80}}{\HLMarginOfError{20.608}}\right)^2=\left(\frac{206.08}{20.608}\right)^2=10^2=\HLObservationCount{100}\text{ couriers}.\]
The response units cancel in the ratio. Thus $\bar x=425$ thousand COP and $n=100$ exactly; no upward rounding is used.
\textbf{Endpoint verification.} The recovered sampling fraction is $100/4000\approx 2.50\%<10\%$. Recompute
\[\begin{aligned}
\HLStandardError{\mathrm{SE}}&=\frac{80}{\sqrt{100}}=\frac{80}{10}=\HLStandardError{8},\qquad \HLMarginOfError{E}=2.576(8)=\HLMarginOfError{20.608},\\
[\bar x-E,\bar x+E]&=[425-20.608,425+20.608]=[\HLIntervalBounds{404.392},\HLIntervalBounds{445.608}].
\end{aligned}\]
This reproduces the supplied endpoints exactly relative to $z^*=2.576$. We are 99\% confident that the cooperative's mean reference-week earnings is between 404.392 and 445.608 thousand COP.
Exact endpoints and the supplied construction cutoff permit exact integer recovery here. Rounded endpoints can make historical $n$ nonunique; see \hyperref[prob:c4-advanced-2]{Advanced Problem 4.2}.
\secInicio[prob:c4-1][sec:general-worked]
\clearpage
\subsubsection*{C4 -- Stage 3: recover $\bar x$ and $\sigma$}\label{sol:c4-1-c}
Each report estimates the cooperative's mean reference-week earnings. Work independently, using the cutoff stated in that report. All response calculations on this page are in thousand COP; $n$ is a count of couriers.
The stated known-population-$\sigma$ procedure identifies the omitted quantity as the population SD used by the analyst. It is not a sample SD $s$.
\par\noindent\textbf{Report 3: midpoint and half-width first.}
\[\begin{aligned}
\HLSampleMean{\bar x}&=\frac{L+U}{2}=\frac{416.84+443.16}{2}=\frac{860}{2}=\HLSampleMean{430},\\
\HLMarginOfError{E}&=\frac{U-L}{2}=\frac{443.16-416.84}{2}=\frac{26.32}{2}=\HLMarginOfError{13.16}.
\end{aligned}\]
\[E=\frac{z^*\sigma}{\sqrt n}\ \Longrightarrow\ E\sqrt n=z^*\sigma\ \Longrightarrow\ \sigma=\frac{E\sqrt n}{z^*}.\]
\[\HLPopulationSD{\sigma}=\frac{\HLMarginOfError{13.16}\sqrt{\HLObservationCount{100}}}{\HLCriticalValue{1.645}}=\frac{13.16(10)}{1.645}=\frac{131.6}{1.645}=\HLPopulationSD{80}\text{ thousand COP}.\]
Thus $\bar x=430$ thousand COP and the analyst used the positive population SD $\sigma=80$ thousand COP.
\textbf{Endpoint verification.} The recovered sampling fraction is $100/4000\approx 2.50\%<10\%$. Recompute
\[\begin{aligned}
\HLStandardError{\mathrm{SE}}&=\frac{80}{\sqrt{100}}=\frac{80}{10}=\HLStandardError{8},\qquad \HLMarginOfError{E}=1.645(8)=\HLMarginOfError{13.16},\\
[\bar x-E,\bar x+E]&=[430-13.16,430+13.16]=[\HLIntervalBounds{416.84},\HLIntervalBounds{443.16}].
\end{aligned}\]
This reproduces the supplied endpoints exactly relative to $z^*=1.645$. We are 90\% confident that the cooperative's mean reference-week earnings is between 416.84 and 443.16 thousand COP.
The recovered SD 80 thousand COP describes individual population values; SE 8 thousand COP describes the sampling spread of their mean. The interval estimates $\mu$, not individual values.
\par\noindent\textbf{Report 4: midpoint and half-width first.}
\[\begin{aligned}
\HLSampleMean{\bar x}&=\frac{L+U}{2}=\frac{405.2+424.8}{2}=\frac{830}{2}=\HLSampleMean{415},\\
\HLMarginOfError{E}&=\frac{U-L}{2}=\frac{424.8-405.2}{2}=\frac{19.6}{2}=\HLMarginOfError{9.8}.
\end{aligned}\]
\[E=\frac{z^*\sigma}{\sqrt n}\ \Longrightarrow\ E\sqrt n=z^*\sigma\ \Longrightarrow\ \sigma=\frac{E\sqrt n}{z^*}.\]
\[\HLPopulationSD{\sigma}=\frac{\HLMarginOfError{9.8}\sqrt{\HLObservationCount{144}}}{\HLCriticalValue{1.96}}=\frac{9.8(12)}{1.96}=\frac{117.6}{1.96}=\HLPopulationSD{60}\text{ thousand COP}.\]
Thus $\bar x=415$ thousand COP and the analyst used the positive population SD $\sigma=60$ thousand COP.
\textbf{Endpoint verification.} The recovered sampling fraction is $144/4000\approx 3.60\%<10\%$. Recompute
\[\begin{aligned}
\HLStandardError{\mathrm{SE}}&=\frac{60}{\sqrt{144}}=\frac{60}{12}=\HLStandardError{5},\qquad \HLMarginOfError{E}=1.96(5)=\HLMarginOfError{9.8},\\
[\bar x-E,\bar x+E]&=[415-9.8,415+9.8]=[\HLIntervalBounds{405.2},\HLIntervalBounds{424.8}].
\end{aligned}\]
This reproduces the supplied endpoints exactly relative to $z^*=1.96$. We are 95\% confident that the cooperative's mean reference-week earnings is between 405.2 and 424.8 thousand COP.
The recovered SD 60 thousand COP describes individual population values; SE 5 thousand COP describes the sampling spread of their mean. The interval estimates $\mu$, not individual values.
\secInicio[prob:c4-1][sec:general-worked]
\clearpage
\subsection{Warehouse charges: construct and recover}\label{prob:c4-2}
An importer studies the quarter's handling charges across 2,400 pallets, one receipt per pallet.
In every case, use a simple random sample of distinct pallets without replacement.
Assume no extreme skew or outliers and an approximately normal sample mean for all the stated
or recovered sizes. The sample fractions are below 10\%, supporting approximate independence;
use the stated known-$\sigma$ $z$ model without a finite-population correction.
The population size is unchanged across the separate scenarios; report-specific means and SDs may differ.
All response values, including SDs and interval endpoints, are in USD per pallet.
\QuestionsTasks[sec:general-worked]{
{4}/{sol:c4-2-a}/{\textbf{Stage 1 -- Construct four intervals.}
Use $\bar x=74$ and known $\sigma=12$ USD per pallet throughout.
For $n_1=36$ and $n_2=144$, construct all four two-sided intervals at
$C_1=95\%$ ($z^*=1.96$) and $C_2=99\%$ ($z^*=2.576$).
Show the four combinations in a comparison table, including margins and widths.
Interpret the estimates and explain both width effects.},
{4}/{sol:c4-2-b}/{\textbf{Stage 2 -- Recover $n$ and $\bar x$.} For EACH independent report below, recover both the sample size and the sample mean. The known-population-$\sigma$ $z$ procedure was used; $n$ and $\bar x$ are omitted.
\newline Report 1: interval $[68.848,79.152]$ USD per pallet; $C=99\%$, construction cutoff $z^*=2.576$, known $\sigma=12$ USD per pallet.
\newline Report 2: interval $[75.5325,80.4675]$ USD per pallet; $C=90\%$, construction cutoff $z^*=1.645$, known $\sigma=15$ USD per pallet.
\newline Find each midpoint and half-width first; show inverse algebra, reconstruct the endpoints and interpret each interval.},
{4}/{sol:c4-2-c}/{\textbf{Stage 3 -- Recover $\bar x$ and $\sigma$.} For EACH further independent report below, recover both the sample mean and the population SD. Each report was constructed using the \textbf{known-population-$\sigma$ $z$ procedure}; the population SD used by the analyst was omitted from the report, as was $\bar x$.
\newline Report 3: interval $[69.06,74.94]$ USD per pallet; $C=95\%$, construction cutoff $z^*=1.96$, sample size $n=64$.
\newline Report 4: interval $[76.136,83.864]$ USD per pallet; $C=99\%$, construction cutoff $z^*=2.576$, sample size $n=144$.
\newline Find each midpoint and half-width first; show inverse algebra, reconstruct the endpoints and interpret the recovered quantities. Explain why the recovered SD is $\sigma$, not $s$.}
}
\clearpage
\subsubsection*{C4 -- Stage 1: four interval constructions}\label{sol:c4-2-a}
The parameter $\HLPopulationMean{\mu}$ is the mean quarterly handling charge across the importer's pallets, in USD per pallet. The sampling and shape assumptions justify the approximate normal model. The known \HLPopulationSD{population SD} is $\HLPopulationSD{\sigma=12}$ USD per pallet; the fixed \HLSampleMean{sample mean} is $\HLSampleMean{\bar x=74}$ USD per pallet.
For each supplied confidence level, $z^*\approx\Phi^{-1}((1+C)/2)$.
Use the supplied cutoffs in the numerical calculations. The formulas are
\[
\HLStandardError{\mathrm{SE}}=\frac{\HLPopulationSD{\sigma}}{\sqrt{\HLObservationCount{n}}},\quad
\HLMarginOfError{E}=\HLCriticalValue{z^*}\HLStandardError{\mathrm{SE}},\quad
[\HLIntervalBounds{L},\HLIntervalBounds{U}]=[\HLSampleMean{\bar x}-\HLMarginOfError{E},\HLSampleMean{\bar x}+\HLMarginOfError{E}],\quad W=2E.
\]
\textbf{$n=36$, $C=95\%$.}
\[\begin{aligned}
\HLStandardError{\mathrm{SE}}&=\frac{\HLPopulationSD{12}}{\sqrt{\HLObservationCount{36}}}=\frac{12}{6}=\HLStandardError{2},\qquad \HLMarginOfError{E}=\HLCriticalValue{1.96}(\HLStandardError{2})=\HLMarginOfError{3.92},\\
[L,U]&=[\HLSampleMean{74}-\HLMarginOfError{3.92},\HLSampleMean{74}+\HLMarginOfError{3.92}]=[\HLIntervalBounds{70.08},\HLIntervalBounds{77.92}],\qquad W=2(3.92)=7.84.
\end{aligned}\]
\textbf{$n=36$, $C=99\%$.}
\[\begin{aligned}
\HLStandardError{\mathrm{SE}}&=\frac{\HLPopulationSD{12}}{\sqrt{\HLObservationCount{36}}}=\frac{12}{6}=\HLStandardError{2},\qquad \HLMarginOfError{E}=\HLCriticalValue{2.576}(\HLStandardError{2})=\HLMarginOfError{5.152},\\
[L,U]&=[\HLSampleMean{74}-\HLMarginOfError{5.152},\HLSampleMean{74}+\HLMarginOfError{5.152}]=[\HLIntervalBounds{68.848},\HLIntervalBounds{79.152}],\qquad W=2(5.152)=10.304.
\end{aligned}\]
\textbf{$n=144$, $C=95\%$.}
\[\begin{aligned}
\HLStandardError{\mathrm{SE}}&=\frac{\HLPopulationSD{12}}{\sqrt{\HLObservationCount{144}}}=\frac{12}{12}=\HLStandardError{1},\qquad \HLMarginOfError{E}=\HLCriticalValue{1.96}(\HLStandardError{1})=\HLMarginOfError{1.96},\\
[L,U]&=[\HLSampleMean{74}-\HLMarginOfError{1.96},\HLSampleMean{74}+\HLMarginOfError{1.96}]=[\HLIntervalBounds{72.04},\HLIntervalBounds{75.96}],\qquad W=2(1.96)=3.92.
\end{aligned}\]
\textbf{$n=144$, $C=99\%$.}
\[\begin{aligned}
\HLStandardError{\mathrm{SE}}&=\frac{\HLPopulationSD{12}}{\sqrt{\HLObservationCount{144}}}=\frac{12}{12}=\HLStandardError{1},\qquad \HLMarginOfError{E}=\HLCriticalValue{2.576}(\HLStandardError{1})=\HLMarginOfError{2.576},\\
[L,U]&=[\HLSampleMean{74}-\HLMarginOfError{2.576},\HLSampleMean{74}+\HLMarginOfError{2.576}]=[\HLIntervalBounds{71.424},\HLIntervalBounds{76.576}],\qquad W=2(2.576)=5.152.
\end{aligned}\]
All SEs, margins, endpoints and widths above and below are in USD per pallet.
\begin{center}
\renewcommand{\arraystretch}{1.25}
\begin{tabular}{@{}rrrrll@{}}\toprule
$\HLObservationCount{n}$&$\HLConfidence{C}$&$\HLCriticalValue{z^*}$&$\HLMarginOfError{E}$&$[\HLIntervalBounds{L},\HLIntervalBounds{U}]$&Width $W$\\\midrule
$\HLObservationCount{36}$&$\HLConfidence{95\%}$&$\HLCriticalValue{1.96}$&$\HLMarginOfError{3.92}$&$[\HLIntervalBounds{70.08},\HLIntervalBounds{77.92}]$&$7.84$\\
$\HLObservationCount{36}$&$\HLConfidence{99\%}$&$\HLCriticalValue{2.576}$&$\HLMarginOfError{5.152}$&$[\HLIntervalBounds{68.848},\HLIntervalBounds{79.152}]$&$10.304$\\
$\HLObservationCount{144}$&$\HLConfidence{95\%}$&$\HLCriticalValue{1.96}$&$\HLMarginOfError{1.96}$&$[\HLIntervalBounds{72.04},\HLIntervalBounds{75.96}]$&$3.92$\\
$\HLObservationCount{144}$&$\HLConfidence{99\%}$&$\HLCriticalValue{2.576}$&$\HLMarginOfError{2.576}$&$[\HLIntervalBounds{71.424},\HLIntervalBounds{76.576}]$&$5.152$\\
\bottomrule\end{tabular}\end{center}
At each row's confidence level, we are 95\% or 99\% confident, respectively, that the mean quarterly handling charge across the importer's pallets lies within that row's bounds.
These intervals estimate a population mean, not a range of individual values.
Holding $\bar x$ and $\sigma$ fixed, increasing $n$ from 36 to 144 multiplies width by
$\sqrt{36/144}=1/2$: quadrupling the sample size halves width at either confidence level.
Increasing confidence from 95\% to 99\% increases width by the factor
$2.576/1.96\approx 1.3143$ at either sample size, since the cutoff increases.
Every interval remains centered at 74 USD per pallet. These are controlled comparisons;
an actual new random sample could have a different sample mean.
\secInicio[prob:c4-2][sec:general-worked]
\clearpage
\subsubsection*{C4 -- Stage 2: recover $n$ and $\bar x$}\label{sol:c4-2-b}
Each report estimates the mean quarterly handling charge across the importer's pallets. Work independently, using the cutoff stated in that report. All response calculations on this page are in USD per pallet; $n$ is a count of pallets.
\par\noindent\textbf{Report 1: midpoint and half-width first.}
\[\begin{aligned}
\HLSampleMean{\bar x}&=\frac{L+U}{2}=\frac{68.848+79.152}{2}=\frac{148}{2}=\HLSampleMean{74},\\
\HLMarginOfError{E}&=\frac{U-L}{2}=\frac{79.152-68.848}{2}=\frac{10.304}{2}=\HLMarginOfError{5.152}.
\end{aligned}\]
\[E=\frac{z^*\sigma}{\sqrt n}\ \Longrightarrow\ E\sqrt n=z^*\sigma\ \Longrightarrow\ \sqrt n=\frac{z^*\sigma}{E}\ \Longrightarrow\ n=\left(\frac{z^*\sigma}{E}\right)^2.\]
\[\HLObservationCount{n}=\left(\frac{\HLCriticalValue{2.576}\cdot\HLPopulationSD{12}}{\HLMarginOfError{5.152}}\right)^2=\left(\frac{30.912}{5.152}\right)^2=6^2=\HLObservationCount{36}\text{ pallets}.\]
The response units cancel in the ratio. Thus $\bar x=74$ USD per pallet and $n=36$ exactly; no upward rounding is used.
\textbf{Endpoint verification.} The recovered sampling fraction is $36/2400\approx 1.50\%<10\%$. Recompute
\[\begin{aligned}
\HLStandardError{\mathrm{SE}}&=\frac{12}{\sqrt{36}}=\frac{12}{6}=\HLStandardError{2},\qquad \HLMarginOfError{E}=2.576(2)=\HLMarginOfError{5.152},\\
[\bar x-E,\bar x+E]&=[74-5.152,74+5.152]=[\HLIntervalBounds{68.848},\HLIntervalBounds{79.152}].
\end{aligned}\]
This reproduces the supplied endpoints exactly relative to $z^*=2.576$. We are 99\% confident that the mean quarterly handling charge across the importer's pallets is between 68.848 and 79.152 USD per pallet.
\par\noindent\textbf{Report 2: midpoint and half-width first.}
\[\begin{aligned}
\HLSampleMean{\bar x}&=\frac{L+U}{2}=\frac{75.5325+80.4675}{2}=\frac{156}{2}=\HLSampleMean{78},\\
\HLMarginOfError{E}&=\frac{U-L}{2}=\frac{80.4675-75.5325}{2}=\frac{4.935}{2}=\HLMarginOfError{2.4675}.
\end{aligned}\]
\[E=\frac{z^*\sigma}{\sqrt n}\ \Longrightarrow\ E\sqrt n=z^*\sigma\ \Longrightarrow\ \sqrt n=\frac{z^*\sigma}{E}\ \Longrightarrow\ n=\left(\frac{z^*\sigma}{E}\right)^2.\]
\[\HLObservationCount{n}=\left(\frac{\HLCriticalValue{1.645}\cdot\HLPopulationSD{15}}{\HLMarginOfError{2.4675}}\right)^2=\left(\frac{24.675}{2.4675}\right)^2=10^2=\HLObservationCount{100}\text{ pallets}.\]
The response units cancel in the ratio. Thus $\bar x=78$ USD per pallet and $n=100$ exactly; no upward rounding is used.
\textbf{Endpoint verification.} The recovered sampling fraction is $100/2400\approx 4.17\%<10\%$. Recompute
\[\begin{aligned}
\HLStandardError{\mathrm{SE}}&=\frac{15}{\sqrt{100}}=\frac{15}{10}=\HLStandardError{1.5},\qquad \HLMarginOfError{E}=1.645(1.5)=\HLMarginOfError{2.4675},\\
[\bar x-E,\bar x+E]&=[78-2.4675,78+2.4675]=[\HLIntervalBounds{75.5325},\HLIntervalBounds{80.4675}].
\end{aligned}\]
This reproduces the supplied endpoints exactly relative to $z^*=1.645$. We are 90\% confident that the mean quarterly handling charge across the importer's pallets is between 75.5325 and 80.4675 USD per pallet.
Exact endpoints and the supplied construction cutoff permit exact integer recovery here. Rounded endpoints can make historical $n$ nonunique; see \hyperref[prob:c4-advanced-2]{Advanced Problem 4.2}.
\secInicio[prob:c4-2][sec:general-worked]
\clearpage
\subsubsection*{C4 -- Stage 3: recover $\bar x$ and $\sigma$}\label{sol:c4-2-c}
Each report estimates the mean quarterly handling charge across the importer's pallets. Work independently, using the cutoff stated in that report. All response calculations on this page are in USD per pallet; $n$ is a count of pallets.
The stated known-population-$\sigma$ procedure identifies the omitted quantity as the population SD used by the analyst. It is not a sample SD $s$.
\par\noindent\textbf{Report 3: midpoint and half-width first.}
\[\begin{aligned}
\HLSampleMean{\bar x}&=\frac{L+U}{2}=\frac{69.06+74.94}{2}=\frac{144}{2}=\HLSampleMean{72},\\
\HLMarginOfError{E}&=\frac{U-L}{2}=\frac{74.94-69.06}{2}=\frac{5.88}{2}=\HLMarginOfError{2.94}.
\end{aligned}\]
\[E=\frac{z^*\sigma}{\sqrt n}\ \Longrightarrow\ E\sqrt n=z^*\sigma\ \Longrightarrow\ \sigma=\frac{E\sqrt n}{z^*}.\]
\[\HLPopulationSD{\sigma}=\frac{\HLMarginOfError{2.94}\sqrt{\HLObservationCount{64}}}{\HLCriticalValue{1.96}}=\frac{2.94(8)}{1.96}=\frac{23.52}{1.96}=\HLPopulationSD{12}\text{ USD per pallet}.\]
Thus $\bar x=72$ USD per pallet and the analyst used the positive population SD $\sigma=12$ USD per pallet.
\textbf{Endpoint verification.} The recovered sampling fraction is $64/2400\approx 2.67\%<10\%$. Recompute
\[\begin{aligned}
\HLStandardError{\mathrm{SE}}&=\frac{12}{\sqrt{64}}=\frac{12}{8}=\HLStandardError{1.5},\qquad \HLMarginOfError{E}=1.96(1.5)=\HLMarginOfError{2.94},\\
[\bar x-E,\bar x+E]&=[72-2.94,72+2.94]=[\HLIntervalBounds{69.06},\HLIntervalBounds{74.94}].
\end{aligned}\]
This reproduces the supplied endpoints exactly relative to $z^*=1.96$. We are 95\% confident that the mean quarterly handling charge across the importer's pallets is between 69.06 and 74.94 USD per pallet.
The recovered SD 12 USD per pallet describes individual population values; SE 1.5 USD per pallet describes the sampling spread of their mean. The interval estimates $\mu$, not individual values.
\par\noindent\textbf{Report 4: midpoint and half-width first.}
\[\begin{aligned}
\HLSampleMean{\bar x}&=\frac{L+U}{2}=\frac{76.136+83.864}{2}=\frac{160}{2}=\HLSampleMean{80},\\
\HLMarginOfError{E}&=\frac{U-L}{2}=\frac{83.864-76.136}{2}=\frac{7.728}{2}=\HLMarginOfError{3.864}.
\end{aligned}\]
\[E=\frac{z^*\sigma}{\sqrt n}\ \Longrightarrow\ E\sqrt n=z^*\sigma\ \Longrightarrow\ \sigma=\frac{E\sqrt n}{z^*}.\]
\[\HLPopulationSD{\sigma}=\frac{\HLMarginOfError{3.864}\sqrt{\HLObservationCount{144}}}{\HLCriticalValue{2.576}}=\frac{3.864(12)}{2.576}=\frac{46.368}{2.576}=\HLPopulationSD{18}\text{ USD per pallet}.\]
Thus $\bar x=80$ USD per pallet and the analyst used the positive population SD $\sigma=18$ USD per pallet.
\textbf{Endpoint verification.} The recovered sampling fraction is $144/2400\approx 6.00\%<10\%$. Recompute
\[\begin{aligned}
\HLStandardError{\mathrm{SE}}&=\frac{18}{\sqrt{144}}=\frac{18}{12}=\HLStandardError{1.5},\qquad \HLMarginOfError{E}=2.576(1.5)=\HLMarginOfError{3.864},\\
[\bar x-E,\bar x+E]&=[80-3.864,80+3.864]=[\HLIntervalBounds{76.136},\HLIntervalBounds{83.864}].
\end{aligned}\]
This reproduces the supplied endpoints exactly relative to $z^*=2.576$. We are 99\% confident that the mean quarterly handling charge across the importer's pallets is between 76.136 and 83.864 USD per pallet.
The recovered SD 18 USD per pallet describes individual population values; SE 1.5 USD per pallet describes the sampling spread of their mean. The interval estimates $\mu$, not individual values.
\secInicio[prob:c4-2][sec:general-worked]
\clearpage
\subsection{Freight charges: construct and recover}\label{prob:c4-3}
A freight insurer studies processing charges across 5,000 shipments, one charge per shipment.
In every case, use a simple random sample of distinct shipments without replacement.
Assume no extreme skew or outliers and an approximately normal sample mean for all the stated
or recovered sizes. The sample fractions are below 10\%, supporting approximate independence;
use the stated known-$\sigma$ $z$ model without a finite-population correction.
The population size is unchanged across the separate scenarios; report-specific means and SDs may differ.
All response values, including SDs and interval endpoints, are in USD per shipment.
\QuestionsTasks[sec:general-worked]{
{4}/{sol:c4-3-a}/{\textbf{Stage 1 -- Construct four intervals.}
Use $\bar x=120$ and known $\sigma=27$ USD per shipment throughout.
For $n_1=81$ and $n_2=324$, construct all four two-sided intervals at
$C_1=90\%$ ($z^*=1.645$) and $C_2=95\%$ ($z^*=1.96$).
Show the four combinations in a comparison table, including margins and widths.
Interpret the estimates and explain both width effects.},
{4}/{sol:c4-3-b}/{\textbf{Stage 2 -- Recover $n$ and $\bar x$.} For EACH independent report below, recover both the sample size and the sample mean. The known-population-$\sigma$ $z$ procedure was used; $n$ and $\bar x$ are omitted.
\newline Report 1: interval $[115.065,124.935]$ USD per shipment; $C=90\%$, construction cutoff $z^*=1.645$, known $\sigma=27$ USD per shipment.
\newline Report 2: interval $[121.1,130.9]$ USD per shipment; $C=95\%$, construction cutoff $z^*=1.96$, known $\sigma=30$ USD per shipment.
\newline Find each midpoint and half-width first; show inverse algebra, reconstruct the endpoints and interpret each interval.},
{4}/{sol:c4-3-c}/{\textbf{Stage 3 -- Recover $\bar x$ and $\sigma$.} For EACH further independent report below, recover both the sample mean and the population SD. Each report was constructed using the \textbf{known-population-$\sigma$ $z$ procedure}; the population SD used by the analyst was omitted from the report, as was $\bar x$.
\newline Report 3: interval $[117.065,126.935]$ USD per shipment; $C=90\%$, construction cutoff $z^*=1.645$, sample size $n=64$.
\newline Report 4: interval $[124.848,135.152]$ USD per shipment; $C=99\%$, construction cutoff $z^*=2.576$, sample size $n=256$.
\newline Find each midpoint and half-width first; show inverse algebra, reconstruct the endpoints and interpret the recovered quantities. Explain why the recovered SD is $\sigma$, not $s$.}
}
\clearpage
\subsubsection*{C4 -- Stage 1: four interval constructions}\label{sol:c4-3-a}
The parameter $\HLPopulationMean{\mu}$ is the mean processing charge across the insurer's shipments, in USD per shipment. The sampling and shape assumptions justify the approximate normal model. The known \HLPopulationSD{population SD} is $\HLPopulationSD{\sigma=27}$ USD per shipment; the fixed \HLSampleMean{sample mean} is $\HLSampleMean{\bar x=120}$ USD per shipment.
For each supplied confidence level, $z^*\approx\Phi^{-1}((1+C)/2)$.
Use the supplied cutoffs in the numerical calculations. The formulas are
\[
\HLStandardError{\mathrm{SE}}=\frac{\HLPopulationSD{\sigma}}{\sqrt{\HLObservationCount{n}}},\quad
\HLMarginOfError{E}=\HLCriticalValue{z^*}\HLStandardError{\mathrm{SE}},\quad
[\HLIntervalBounds{L},\HLIntervalBounds{U}]=[\HLSampleMean{\bar x}-\HLMarginOfError{E},\HLSampleMean{\bar x}+\HLMarginOfError{E}],\quad W=2E.
\]
\textbf{$n=81$, $C=90\%$.}
\[\begin{aligned}
\HLStandardError{\mathrm{SE}}&=\frac{\HLPopulationSD{27}}{\sqrt{\HLObservationCount{81}}}=\frac{27}{9}=\HLStandardError{3},\qquad \HLMarginOfError{E}=\HLCriticalValue{1.645}(\HLStandardError{3})=\HLMarginOfError{4.935},\\
[L,U]&=[\HLSampleMean{120}-\HLMarginOfError{4.935},\HLSampleMean{120}+\HLMarginOfError{4.935}]=[\HLIntervalBounds{115.065},\HLIntervalBounds{124.935}],\qquad W=2(4.935)=9.87.
\end{aligned}\]
\textbf{$n=81$, $C=95\%$.}
\[\begin{aligned}
\HLStandardError{\mathrm{SE}}&=\frac{\HLPopulationSD{27}}{\sqrt{\HLObservationCount{81}}}=\frac{27}{9}=\HLStandardError{3},\qquad \HLMarginOfError{E}=\HLCriticalValue{1.96}(\HLStandardError{3})=\HLMarginOfError{5.88},\\
[L,U]&=[\HLSampleMean{120}-\HLMarginOfError{5.88},\HLSampleMean{120}+\HLMarginOfError{5.88}]=[\HLIntervalBounds{114.12},\HLIntervalBounds{125.88}],\qquad W=2(5.88)=11.76.
\end{aligned}\]
\textbf{$n=324$, $C=90\%$.}
\[\begin{aligned}
\HLStandardError{\mathrm{SE}}&=\frac{\HLPopulationSD{27}}{\sqrt{\HLObservationCount{324}}}=\frac{27}{18}=\HLStandardError{1.5},\qquad \HLMarginOfError{E}=\HLCriticalValue{1.645}(\HLStandardError{1.5})=\HLMarginOfError{2.4675},\\
[L,U]&=[\HLSampleMean{120}-\HLMarginOfError{2.4675},\HLSampleMean{120}+\HLMarginOfError{2.4675}]=[\HLIntervalBounds{117.5325},\HLIntervalBounds{122.4675}],\qquad W=2(2.4675)=4.935.
\end{aligned}\]
\textbf{$n=324$, $C=95\%$.}
\[\begin{aligned}
\HLStandardError{\mathrm{SE}}&=\frac{\HLPopulationSD{27}}{\sqrt{\HLObservationCount{324}}}=\frac{27}{18}=\HLStandardError{1.5},\qquad \HLMarginOfError{E}=\HLCriticalValue{1.96}(\HLStandardError{1.5})=\HLMarginOfError{2.94},\\
[L,U]&=[\HLSampleMean{120}-\HLMarginOfError{2.94},\HLSampleMean{120}+\HLMarginOfError{2.94}]=[\HLIntervalBounds{117.06},\HLIntervalBounds{122.94}],\qquad W=2(2.94)=5.88.
\end{aligned}\]
All SEs, margins, endpoints and widths above and below are in USD per shipment.
\begin{center}
\renewcommand{\arraystretch}{1.25}
\begin{tabular}{@{}rrrrll@{}}\toprule
$\HLObservationCount{n}$&$\HLConfidence{C}$&$\HLCriticalValue{z^*}$&$\HLMarginOfError{E}$&$[\HLIntervalBounds{L},\HLIntervalBounds{U}]$&Width $W$\\\midrule
$\HLObservationCount{81}$&$\HLConfidence{90\%}$&$\HLCriticalValue{1.645}$&$\HLMarginOfError{4.935}$&$[\HLIntervalBounds{115.065},\HLIntervalBounds{124.935}]$&$9.87$\\
$\HLObservationCount{81}$&$\HLConfidence{95\%}$&$\HLCriticalValue{1.96}$&$\HLMarginOfError{5.88}$&$[\HLIntervalBounds{114.12},\HLIntervalBounds{125.88}]$&$11.76$\\
$\HLObservationCount{324}$&$\HLConfidence{90\%}$&$\HLCriticalValue{1.645}$&$\HLMarginOfError{2.4675}$&$[\HLIntervalBounds{117.5325},\HLIntervalBounds{122.4675}]$&$4.935$\\
$\HLObservationCount{324}$&$\HLConfidence{95\%}$&$\HLCriticalValue{1.96}$&$\HLMarginOfError{2.94}$&$[\HLIntervalBounds{117.06},\HLIntervalBounds{122.94}]$&$5.88$\\
\bottomrule\end{tabular}\end{center}
At each row's confidence level, we are 90\% or 95\% confident, respectively, that the mean processing charge across the insurer's shipments lies within that row's bounds.
These intervals estimate a population mean, not a range of individual values.
Holding $\bar x$ and $\sigma$ fixed, increasing $n$ from 81 to 324 multiplies width by
$\sqrt{81/324}=1/2$: quadrupling the sample size halves width at either confidence level.
Increasing confidence from 90\% to 95\% increases width by the factor
$1.96/1.645\approx 1.1915$ at either sample size, since the cutoff increases.
Every interval remains centered at 120 USD per shipment. These are controlled comparisons;
an actual new random sample could have a different sample mean.
\secInicio[prob:c4-3][sec:general-worked]
\clearpage
\subsubsection*{C4 -- Stage 2: recover $n$ and $\bar x$}\label{sol:c4-3-b}
Each report estimates the mean processing charge across the insurer's shipments. Work independently, using the cutoff stated in that report. All response calculations on this page are in USD per shipment; $n$ is a count of shipments.
\par\noindent\textbf{Report 1: midpoint and half-width first.}
\[\begin{aligned}
\HLSampleMean{\bar x}&=\frac{L+U}{2}=\frac{115.065+124.935}{2}=\frac{240}{2}=\HLSampleMean{120},\\
\HLMarginOfError{E}&=\frac{U-L}{2}=\frac{124.935-115.065}{2}=\frac{9.87}{2}=\HLMarginOfError{4.935}.
\end{aligned}\]
\[E=\frac{z^*\sigma}{\sqrt n}\ \Longrightarrow\ E\sqrt n=z^*\sigma\ \Longrightarrow\ \sqrt n=\frac{z^*\sigma}{E}\ \Longrightarrow\ n=\left(\frac{z^*\sigma}{E}\right)^2.\]
\[\HLObservationCount{n}=\left(\frac{\HLCriticalValue{1.645}\cdot\HLPopulationSD{27}}{\HLMarginOfError{4.935}}\right)^2=\left(\frac{44.415}{4.935}\right)^2=9^2=\HLObservationCount{81}\text{ shipments}.\]
The response units cancel in the ratio. Thus $\bar x=120$ USD per shipment and $n=81$ exactly; no upward rounding is used.
\textbf{Endpoint verification.} The recovered sampling fraction is $81/5000\approx 1.62\%<10\%$. Recompute
\[\begin{aligned}
\HLStandardError{\mathrm{SE}}&=\frac{27}{\sqrt{81}}=\frac{27}{9}=\HLStandardError{3},\qquad \HLMarginOfError{E}=1.645(3)=\HLMarginOfError{4.935},\\
[\bar x-E,\bar x+E]&=[120-4.935,120+4.935]=[\HLIntervalBounds{115.065},\HLIntervalBounds{124.935}].
\end{aligned}\]
This reproduces the supplied endpoints exactly relative to $z^*=1.645$. We are 90\% confident that the mean processing charge across the insurer's shipments is between 115.065 and 124.935 USD per shipment.
\par\noindent\textbf{Report 2: midpoint and half-width first.}
\[\begin{aligned}
\HLSampleMean{\bar x}&=\frac{L+U}{2}=\frac{121.1+130.9}{2}=\frac{252}{2}=\HLSampleMean{126},\\
\HLMarginOfError{E}&=\frac{U-L}{2}=\frac{130.9-121.1}{2}=\frac{9.8}{2}=\HLMarginOfError{4.9}.
\end{aligned}\]
\[E=\frac{z^*\sigma}{\sqrt n}\ \Longrightarrow\ E\sqrt n=z^*\sigma\ \Longrightarrow\ \sqrt n=\frac{z^*\sigma}{E}\ \Longrightarrow\ n=\left(\frac{z^*\sigma}{E}\right)^2.\]
\[\HLObservationCount{n}=\left(\frac{\HLCriticalValue{1.96}\cdot\HLPopulationSD{30}}{\HLMarginOfError{4.9}}\right)^2=\left(\frac{58.8}{4.9}\right)^2=12^2=\HLObservationCount{144}\text{ shipments}.\]
The response units cancel in the ratio. Thus $\bar x=126$ USD per shipment and $n=144$ exactly; no upward rounding is used.
\textbf{Endpoint verification.} The recovered sampling fraction is $144/5000\approx 2.88\%<10\%$. Recompute
\[\begin{aligned}
\HLStandardError{\mathrm{SE}}&=\frac{30}{\sqrt{144}}=\frac{30}{12}=\HLStandardError{2.5},\qquad \HLMarginOfError{E}=1.96(2.5)=\HLMarginOfError{4.9},\\
[\bar x-E,\bar x+E]&=[126-4.9,126+4.9]=[\HLIntervalBounds{121.1},\HLIntervalBounds{130.9}].
\end{aligned}\]
This reproduces the supplied endpoints exactly relative to $z^*=1.96$. We are 95\% confident that the mean processing charge across the insurer's shipments is between 121.1 and 130.9 USD per shipment.
Exact endpoints and the supplied construction cutoff permit exact integer recovery here. Rounded endpoints can make historical $n$ nonunique; see \hyperref[prob:c4-advanced-2]{Advanced Problem 4.2}.
\secInicio[prob:c4-3][sec:general-worked]
\clearpage
\subsubsection*{C4 -- Stage 3: recover $\bar x$ and $\sigma$}\label{sol:c4-3-c}
Each report estimates the mean processing charge across the insurer's shipments. Work independently, using the cutoff stated in that report. All response calculations on this page are in USD per shipment; $n$ is a count of shipments.
The stated known-population-$\sigma$ procedure identifies the omitted quantity as the population SD used by the analyst. It is not a sample SD $s$.
\par\noindent\textbf{Report 3: midpoint and half-width first.}
\[\begin{aligned}
\HLSampleMean{\bar x}&=\frac{L+U}{2}=\frac{117.065+126.935}{2}=\frac{244}{2}=\HLSampleMean{122},\\
\HLMarginOfError{E}&=\frac{U-L}{2}=\frac{126.935-117.065}{2}=\frac{9.87}{2}=\HLMarginOfError{4.935}.
\end{aligned}\]
\[E=\frac{z^*\sigma}{\sqrt n}\ \Longrightarrow\ E\sqrt n=z^*\sigma\ \Longrightarrow\ \sigma=\frac{E\sqrt n}{z^*}.\]
\[\HLPopulationSD{\sigma}=\frac{\HLMarginOfError{4.935}\sqrt{\HLObservationCount{64}}}{\HLCriticalValue{1.645}}=\frac{4.935(8)}{1.645}=\frac{39.48}{1.645}=\HLPopulationSD{24}\text{ USD per shipment}.\]
Thus $\bar x=122$ USD per shipment and the analyst used the positive population SD $\sigma=24$ USD per shipment.
\textbf{Endpoint verification.} The recovered sampling fraction is $64/5000\approx 1.28\%<10\%$. Recompute
\[\begin{aligned}
\HLStandardError{\mathrm{SE}}&=\frac{24}{\sqrt{64}}=\frac{24}{8}=\HLStandardError{3},\qquad \HLMarginOfError{E}=1.645(3)=\HLMarginOfError{4.935},\\
[\bar x-E,\bar x+E]&=[122-4.935,122+4.935]=[\HLIntervalBounds{117.065},\HLIntervalBounds{126.935}].
\end{aligned}\]
This reproduces the supplied endpoints exactly relative to $z^*=1.645$. We are 90\% confident that the mean processing charge across the insurer's shipments is between 117.065 and 126.935 USD per shipment.
The recovered SD 24 USD per shipment describes individual population values; SE 3 USD per shipment describes the sampling spread of their mean. The interval estimates $\mu$, not individual values.
\par\noindent\textbf{Report 4: midpoint and half-width first.}
\[\begin{aligned}
\HLSampleMean{\bar x}&=\frac{L+U}{2}=\frac{124.848+135.152}{2}=\frac{260}{2}=\HLSampleMean{130},\\
\HLMarginOfError{E}&=\frac{U-L}{2}=\frac{135.152-124.848}{2}=\frac{10.304}{2}=\HLMarginOfError{5.152}.
\end{aligned}\]
\[E=\frac{z^*\sigma}{\sqrt n}\ \Longrightarrow\ E\sqrt n=z^*\sigma\ \Longrightarrow\ \sigma=\frac{E\sqrt n}{z^*}.\]
\[\HLPopulationSD{\sigma}=\frac{\HLMarginOfError{5.152}\sqrt{\HLObservationCount{256}}}{\HLCriticalValue{2.576}}=\frac{5.152(16)}{2.576}=\frac{82.432}{2.576}=\HLPopulationSD{32}\text{ USD per shipment}.\]
Thus $\bar x=130$ USD per shipment and the analyst used the positive population SD $\sigma=32$ USD per shipment.
\textbf{Endpoint verification.} The recovered sampling fraction is $256/5000\approx 5.12\%<10\%$. Recompute
\[\begin{aligned}
\HLStandardError{\mathrm{SE}}&=\frac{32}{\sqrt{256}}=\frac{32}{16}=\HLStandardError{2},\qquad \HLMarginOfError{E}=2.576(2)=\HLMarginOfError{5.152},\\
[\bar x-E,\bar x+E]&=[130-5.152,130+5.152]=[\HLIntervalBounds{124.848},\HLIntervalBounds{135.152}].
\end{aligned}\]
This reproduces the supplied endpoints exactly relative to $z^*=2.576$. We are 99\% confident that the mean processing charge across the insurer's shipments is between 124.848 and 135.152 USD per shipment.
The recovered SD 32 USD per shipment describes individual population values; SE 2 USD per shipment describes the sampling spread of their mean. The interval estimates $\mu$, not individual values.
\secInicio[prob:c4-3][sec:general-worked]

\clearpage
\Needspace{20\baselineskip}
\section{Advanced Problems}\label{sec:advanced}
\SectionProblemsTableOfContents{
\SectionProblemLink{prob:c4-advanced-1}{Compare two warehouse audits (C4)}\\
\SectionProblemLink{prob:c4-advanced-2}{Rounded courier reports and nonunique recovery (C4)}
}
\subsection{Compare two warehouse audits}\label{prob:c4-advanced-1}
Two warehouses each have 5,000 pallet receipts. Each audit uses a simple random sample of distinct pallets without replacement, one receipt per pallet. Charges are approximately bell-shaped. Both reports use the known-population-$\sigma$ $z$ procedure at 95\%, with construction cutoff $z^*=1.96$ and exact endpoints relative to that cutoff.
\begin{center}
\begin{tabular}{@{}lccc@{}}\toprule
Audit&Interval (USD per pallet)&Known population SD&Sample size\\\midrule
A&$[56.08,63.92]$&16 USD&omitted\\
B&$[57.06,62.94]$&omitted&64\\\bottomrule
\end{tabular}
\end{center}
\QuestionsTasks[sec:advanced]{
{4}/{sol:c4-advanced-1}/{Recover both sample means and the missing quantity in each audit. Compare widths using both sample size and population SD. For audit B construct a 99\% interval with $z^*=2.576$; is it narrower than A's 95\% interval? Can equal sample means establish equal population means?}
}
\subsubsection*{C4 -- Coordinate all three problem types}\label{sol:c4-advanced-1}
Let $\mu_A,\mu_B$ be the warehouses' population mean charges. Random sampling and normal charge shapes support the $z$ model; check sampling fractions below.

For A, $\HLSampleMean{\bar x_A}=(56.08+63.92)/2=\HLSampleMean{60}$ USD and $\HLMarginOfError{E_A}=(63.92-56.08)/2=\HLMarginOfError{3.92}$ USD.
From $E_A\sqrt{n_A}=z^*\sigma_A$, obtain
\[
n_A=\left(\frac{z^*\sigma_A}{E_A}\right)^2
=\left(\frac{1.96(16)}{3.92}\right)^2
=\left(\frac{31.36}{3.92}\right)^2=8^2=64.
\]
For B, $\bar x_B=(57.06+62.94)/2=60$ USD and $E_B=(62.94-57.06)/2=2.94$ USD. Rearranging $E_B=z^*\sigma_B/\sqrt{n_B}$ gives
\[
\sigma_B=\frac{E_B\sqrt{n_B}}{z^*}
=\frac{2.94(8)}{1.96}=\frac{23.52}{1.96}=12\text{ USD}.
\]
Both sampling fractions are $64/5{,}000=1.28\%$. Check: A has $\mathrm{SE}_A=16/8=2$ USD and $60\pm1.96(2)=[56.08,63.92]$ USD; B has $\mathrm{SE}_B=12/8=1.5$ USD and $60\pm1.96(1.5)=[57.06,62.94]$ USD.

Widths are $W_A=7.84$ and $W_B=5.88$ USD. Equal sample sizes and confidence levels give $W_B/W_A=\sigma_B/\sigma_A=12/16=0.75$: B's smaller population spread explains its narrower interval.

For B at 99\%, $E=2.576(12)/8=3.864$ USD, giving $[56.136,63.864]$ USD and width $7.728$ USD. It is wider than B's 95\% interval, yet narrower than A's 95\% width of 7.84 USD. Confidence alone does not determine width across populations. Equal sample means of 60 USD do not prove $\mu_A=\mu_B$; each interval estimates its own population mean.
\secInicio[prob:c4-advanced-1][sec:advanced]

\clearpage
\subsection{Rounded courier reports and nonunique recovery}\label{prob:c4-advanced-2}
A courier report displays a 95\% interval $[390,430]$ thousand COP, with each endpoint rounded to the nearest whole thousand COP. It used the known-population-$\sigma$ $z$ procedure with construction cutoff $z^*=1.96$ and known population SD 80 thousand COP. A simple random sample of distinct couriers was drawn from 2,400; earnings have no extreme skew or outliers, and the candidate sample sizes below support approximately normal sample means. For this comparison, suppose each candidate sample has mean 410 thousand COP.
\QuestionsTasks[sec:advanced]{
{4}/{sol:c4-advanced-2}/{Invert the displayed endpoints as though exact. Then construct unrounded intervals for $n=60$ and $n=64$ and round them as the report did. Can the report uniquely recover $n$ or the exact original sample mean? Explain why rounding the inverse result upward is unjustified.}
}
\subsubsection*{C4 -- Test what the reported precision supports}\label{sol:c4-advanced-2}
The parameter is the roster's mean reference-week earnings. The candidate fractions are $60/2{,}400=2.5\%$ and $64/2{,}400\approx2.7\%$; random sampling, small fractions and the stated shape support the approximate known-$\sigma$ normal model.

If the displayed endpoints were exact, the midpoint would be $(390+430)/2=410$ and $E=(430-390)/2=20$ thousand COP. Rearranging $E\sqrt n=z^*\sigma$ gives
\[
n=\left(\frac{z^*\sigma}{E}\right)^2
=\left(\frac{1.96(80)}{20}\right)^2
=7.84^2=61.4656.
\]
This is not a positive integer. Here the endpoints are rounded, so it does not identify the original sample size. Test two possibilities with the forward model:
\[
\begin{aligned}
n=60:\quad &\mathrm{SE}=80/\sqrt{60}\approx10.32796,\\
&E=1.96(80)/\sqrt{60}\approx20.24279,\\
&[\bar x-E,\bar x+E]\approx[389.75721,430.24279];\\[3pt]
n=64:\quad &\mathrm{SE}=80/8=10,\qquad E=1.96(10)=19.6,\\
&[\bar x-E,\bar x+E]=[390.4,429.6].
\end{aligned}
\]
All monetary values are in thousand COP. Both intervals round to $[390,430]$, so the displayed report cannot distinguish $n=60$ from $n=64$. Larger $n$ reduces the unrounded width here, but coarse reporting hides the difference.

The exact sample mean is not uniquely determined either. Even for $n=64$, a sample mean of 410.05 gives $[390.45,429.65]$ thousand COP, which also rounds to $[390,430]$. The midpoint of displayed endpoints is only an approximate reconstruction.

Taking $\lceil61.4656\rceil=62$ would answer a different question: the minimum planned $n$ to achieve $E\leq20$ thousand COP at this confidence and known SD. It does not recover the historical sample size. Request greater endpoint precision or the original sample count to distinguish the possibilities. Interpret the report as an approximate interval for the population mean, never as a range for individual courier earnings.
\secInicio[prob:c4-advanced-2][sec:advanced]
\end{document}
